\section{新钱、老钱与国家民族机会}

前面讲公司进入主机厂体系，本章进入更宏观的全球化和产业叙事。李一帆谈到新钱和老钱，也谈到国家、民族和产业机会。硬科技创业与消费互联网不同，它常常绑定国家制造能力、产业链完整性、全球客户沟通和资本文化差异。

\begin{figure}[H]
\centering
\includegraphics[width=0.94\textwidth]{new-money-old-money.png}
\caption{新钱与老钱：全球化沟通里，不同资本和产业文化有不同信任语言。自制概念图，依据 02:38:34--03:02:16 对谈内容整理。}
\end{figure}

\begin{knowledgebox}{读图：全球化不是只翻译语言}
老钱重视历史、秩序和稳健，新钱重视速度、增长和技术。中国硬科技公司全球化时，不只要解释产品，还要解释自己可信、长期、可交付。
\end{knowledgebox}

\subsection{国家、民族与产业机会}

全球化之后，本节回到更大的产业叙事。开头和结尾都出现“国家、民族、机会”的表达。这里不是抽象口号，而是产业现实：新能源汽车、传感器、芯片、供应链、制造和全球市场共同决定硬科技公司是否有机会。某些机会不是单个创业者创造，而是时代和产业给的。

\begin{figure}[H]
\centering
\includegraphics[width=0.92\textwidth]{national-industrial-opportunity.png}
\caption{国家、民族与产业机会：硬科技创业常常来自时代、产业链和国家能力的叠加。自制概念图，依据 00:00:04--00:00:59 与 02:38:34--03:02:16 对谈内容整理。}
\end{figure}

\begin{warningbox}{宏大叙事不能替代产品能力}
国家和民族机会给的是产业窗口，不是免费胜利。企业仍然要靠技术、成本、质量、客户和组织能力兑现窗口。
\end{warningbox}

\begin{knowledgebox}{硬科技公司的最终考试}
\begin{tabularx}{\textwidth}{p{0.22\textwidth}p{0.34\textwidth}X}
\toprule
考试 & 问题 & 通过标准 \\
\midrule
技术考试 & 性能是否真的领先？ & 客观测试、客户场景、长期可靠性。 \\
成本考试 & 能否持续降本？ & BOM、良率、供应链和规模化。 \\
客户考试 & 主机厂是否愿意采用？ & 认证、定点、量产和复购。 \\
资本考试 & 能否穿越周期？ & 融资、现金流、上市和市场信任。 \\
全球化考试 & 能否跨文化交付？ & 海外客户、合规、品牌和服务网络。 \\
\bottomrule
\end{tabularx}
\end{knowledgebox}

\begin{knowledgebox}{全球化沟通清单}
\begin{tabularx}{\textwidth}{p{0.22\textwidth}p{0.34\textwidth}X}
\toprule
维度 & 要证明什么 & 典型材料 \\
\midrule
技术可信 & 性能和可靠性不是宣传口径 & 测试报告、客户案例、车规认证。 \\
商业可信 & 公司能长期供货和服务 & 财务、产能、售后网络。 \\
文化可信 & 能理解不同客户和资本文化 & 本地团队、沟通方式、长期承诺。 \\
合规可信 & 符合当地规则和供应链要求 & 合规审查、数据和安全流程。 \\
\bottomrule
\end{tabularx}
\end{knowledgebox}

\subsection{本章小结}

硬科技创业既是公司故事，也是产业和国家能力故事。理解禾赛，要同时看技术细节和时代窗口。

